% TOP_PROBLEM: {"id": 2815, "title": "Kirby Problem 3.17", "category_id": 7, "category": {"id": 7, "name": "topology", "display_name": "Topology", "description": "Properties preserved under continuous deformations.", "slug": "topology", "order_index": 7, "created_at": "2026-07-31T15:26:25.671Z"}, "status": "open", "problem_number": "KP-3.17", "set_id": 11}
% FIRST_POSTED: 2026-10-01
% ENTRY_KIND: statement_only
% UnsolvedMath ID 2815; source catalogue code KP-3.17
% !TeX program = lualatex
% Definition and sources only; this file is not an LLM solution attempt.
\documentclass[11pt,a4paper]{article}
\usepackage[margin=25mm]{geometry}
\usepackage{fontspec}
\setmainfont{lmroman10-regular.otf}[BoldFont=lmroman10-bold.otf,
 ItalicFont=lmroman10-italic.otf,BoldItalicFont=lmroman10-bolditalic.otf]
\usepackage{amsmath,amssymb,mathtools}
\usepackage{unicode-math}
\setmathfont{latinmodern-math.otf}
\AtBeginDocument{\let\setminus\smallsetminus}
\usepackage{xurl,hyperref}
\hypersetup{colorlinks=true,urlcolor=blue}
\setcounter{secnumdepth}{0}
\setlength{\parindent}{0pt}
\setlength{\parskip}{5pt}
\setlength{\emergencystretch}{3em}
\begin{document}
\section*{Kirby Problem 3.17}\label{2815}
{\small UnsolvedMath ID 2815 · KP-3.17 · Topology · Kirby's Problems in Low-Dimensional Topology}
\subsection{Definitions and mathematical statement}
For a complete finite-volume hyperbolic $n$-manifold $M$, $n\geq2$, its
length spectrum is the multiset of lengths of closed geodesics. One may first
list primitive unoriented closed geodesics, counting distinct geometric
geodesics of the same length separately; including every positive iterate
then gives the full spectrum. Use the same convention for both manifolds.
Equality of the full spectra is equivalent to equality of the primitive
multisets after the contributions of iterates are removed.

Two hyperbolic manifolds are commensurable if there exists a hyperbolic
manifold $N$ with finite-sheeted local-isometry covering maps to both.
This is weaker than being isometric, since distinct finite covers of one
manifold may have identical length spectra.

Suppose $M_1$ and $M_2$ are complete finite-volume hyperbolic manifolds
of the same dimension and have equal length spectra \emph{with multiplicity}.
Must they be commensurable? The lengths are real translation lengths of
geodesics; no extra normal-bundle holonomy angles are supplied. Agreement
up to an arbitrary but finite cutoff is insufficient, and agreement only
as sets after discarding multiplicities is a different hypothesis. The
assertion ranges over both arithmetic and nonarithmetic manifolds.
\subsection{Short English statement}
If two finite-volume hyperbolic manifolds have exactly the same closed-geodesic lengths with multiplicity, must they have a common finite cover?
\subsection{Sources}
\begin{itemize}
\item[S1] ULAM.AI and the UnsolvedMath Contributors, supplied problems.json, record 2815 (KP-3.17), Kirby's Problems in Low-Dimensional Topology. Catalogue identity and source transcription; CC BY 4.0.
\url{https://huggingface.co/datasets/ulamai/UnsolvedMath}.
\item[S2] K3: A New Problem List in Low-Dimensional Topology, Problem 3.17. Expanded formulation of the supplied catalogue statement.
\url{https://math.berkeley.edu/sites/default/files/surv-295-ruberman-watermarked-author-pdf.pdf}.
\end{itemize}
\end{document}
